\chapter{The Ecosystem Map}
\label{ch:ecosystem-map}

CNA does not sit alone in the \texttt{openeggbert} GitHub organization. It sits at the
center of a small, deliberately layered set of sibling repositories, most of which expect to
be checked out as literal sibling directories on disk --- \texttt{../sharp-runtime},
\texttt{../easy-gl}, \texttt{../cna-samples}, and so on --- because CNake's own build system
looks for them there. This chapter is the map; later chapters (Part~VI for
\texttt{sharp-runtime}, Part~VII for \texttt{easy-gl} and \texttt{free-direct}) are the
territory.

\section{The dependency graph}

\begin{center}
\begin{tabular}{lll}
\toprule
\textbf{Repository} & \textbf{Role} & \textbf{Consumed by} \\
\midrule
\texttt{sharp-runtime} & C\texttt{++} subset of the .NET BCL & CNA (foundation layer) \\
\texttt{easy-gl} & Toolkit-independent OpenGL/ES wrapper & CNA's \texttt{EASYGL} backend \\
\texttt{free-direct} & DirectDraw/DirectSound/DirectPlay subset & CNA's \texttt{DX3} backend \\
\texttt{xna4-spec} & Machine-readable XNA 4.0 API spec (XML) & Compatibility audits of CNA \\
\texttt{cna-samples} & Ports of official XNA 4.0 samples & Exercises CNA end-to-end \\
\texttt{cna-examples} & In-app demo catalog for CNA & Exercises CNA end-to-end \\
\bottomrule
\end{tabular}
\end{center}

Every one of these is a genuinely independent, separately buildable, separately tested
project --- none of them is a vendored copy or a submodule of CNA. \texttt{sharp-runtime} in
particular describes itself as serving ``as a foundation for higher-level frameworks (e.g.
CNA)'' and lists CNA as its own downstream consumer, not the other way around; you could, in
principle, use \texttt{sharp-runtime} or \texttt{easy-gl} in a project that has nothing to do
with XNA at all.

\section{sharp-runtime: the floor CNA stands on}
\label{sec:sharp-runtime-intro}

XNA's public API vocabulary assumes .NET underneath it: \texttt{byte}, \texttt{TimeSpan},
events, \texttt{IDisposable}, \texttt{System.Object.GetType()}. A C\texttt{++} reimplementation
of XNA has two choices --- invent ad-hoc C\texttt{++} substitutes for all of that inline,
scattered through the XNA porting work itself, or build the .NET-shaped foundation once, as
its own project, and depend on it. CNA's own \texttt{CLAUDE.md} porting rules choose the
second option explicitly: ``If CNA needs anything from .NET that is not yet in
\texttt{sharp-runtime} --- a type alias, a class, an interface, an exception type, a utility
--- add it to \texttt{sharp-runtime} first, then use it from CNA. Do not inline .NET concepts
directly into CNA headers as raw C\texttt{++} types or ad-hoc workarounds.''

\texttt{sharp-runtime} is explicit that it is \emph{not} attempting a CLR: no reflection, no
GC, no delegates beyond \texttt{std::function}, no serialization infrastructure, no P/Invoke,
no full cryptography/TLS stack. What it does provide is documented as a ``pragmatic subset
designed for use in native C\texttt{++} applications,'' tracked class-by-class in a
git-ignored local SQLite database (\texttt{plan.sqlite3}) with an honest, small state machine
per .NET type: \texttt{todo}, \texttt{ported}, \texttt{ignore} / \texttt{ignored}, or
\texttt{tobedecided} for genuinely ambiguous cases that need a human architecture call rather
than a guess. Chapter~\ref{ch:sharp-runtime-overview} and Chapter~\ref{ch:sharp-runtime-namespaces}
cover \texttt{sharp-runtime}'s own architecture, its \texttt{System::*} namespace layout, and its
own roughly 10\,900-test suite in depth; Chapter~\ref{ch:parity-philosophy} covers the specific
question of how strictly CNA insists on that .NET shape versus where it deliberately diverges.

\section{easy-gl: OpenGL without a windowing opinion}

The \texttt{EASYGL} graphics backend --- CNA's most mature backend by its own account, with
full pixel-verified 2D and 3D coverage --- is built on \texttt{easy-gl}, a small,
toolkit-independent C\texttt{++}20 wrapper over OpenGL and OpenGLES. Its defining
architectural choice is what it refuses to own: no window creation, no GL context creation,
no event loop. The host application (CNA, in this case) owns the window, creates and
activates the GL context, and hands \texttt{easy-gl} a \texttt{GetProcAddress}-style loader
callback; \texttt{easy-gl} owns everything from there down --- shader compilation and
linking, buffers, vertex arrays, textures, draw calls, and capability queries that let calling
code gate OpenGL-only features (like polygon-mode fill flags) away from OpenGLES targets at
runtime via \texttt{device.supports(...)} / \texttt{device.require(...)}. As of the snapshot this
book draws from, \texttt{easy-gl}'s own README describes its scope as still a deliberately
compact, evolving vertical slice --- initialization, capability detection, and basic draw flow,
exercised by its own \texttt{hello-triangle-sdl} example and two dedicated smoke-test suites
(\texttt{easy-gl-smoke-tests}, \texttt{easy-gl-resource-smoke-tests}) --- with CNA's own
\texttt{EASYGL} backend having grown well beyond that slice on top of it, per
Chapter~\ref{ch:easygl-backend}. Chapter~\ref{ch:easygl-deep-dive} covers \texttt{easy-gl}'s
public API surface (\texttt{easygl::Device}, \texttt{easygl::Shader}, \texttt{easygl::Program},
\texttt{easygl::Buffer}, \texttt{easygl::VertexArray}, \texttt{easygl::Texture},
\texttt{easygl::Feature}) in full.

\section{free-direct: a narrow DirectX 3 for real legacy games}

Where \texttt{easy-gl} is a general-purpose OpenGL wrapper, \texttt{free-direct} is the
opposite kind of project: a deliberately \emph{narrow}, game-driven reimplementation of a
subset of DirectX~3 (2D) --- DirectDraw, DirectSound, and DirectPlay --- built on SDL3
internally, whose explicit goal is not DirectX compatibility in general but running two
specific legacy Windows games (\emph{Speedy Blupi} and \emph{Planet Blupi}) without their
original OS dependencies. Its own README states the non-goals as plainly as the goals: no
full DirectX~3 compatibility, no hardware-accurate emulation, no Direct3D pipeline, and no
expansion of APIs the target games do not actually call. This project backs CNA's
\texttt{DX3} graphics backend (Chapter~\ref{ch:dx3-freedirect-backend}) and is covered as a
project in its own right, including its real (not stubbed) ENet-backed DirectPlay networking,
in Chapter~\ref{ch:freedirect-deep-dive} and in the Speedy Blupi/Planet Blupi porting case study
in Chapter~\ref{ch:blupi-case-study}. \texttt{free-direct}'s own test suite is deliberately split
by transport: the default \texttt{ctest} run (no label filter) exercises
\texttt{LoopbackDirectPlayTransport}'s synchronous semantics and is not expected to pass
\texttt{directplay\_tests} unmodified against a real, asynchronous ENet build (1 of 8 tests fails
that way, by design, not as a regression); a separate \texttt{-L enet} label opts into the
ENet-backed transport tests specifically, verified 1/1 passing.

\section{xna4-spec: a ground truth to audit against}

A recurring failure mode for any reimplementation project is drift between what the README
claims and what the code actually does --- and a subtler one is drift between what the
README claims and what the \emph{original} API actually specified, from memory rather than
from a real reference. \texttt{xna4-spec} exists to remove the second failure mode. It is a
complete, machine-readable XML transcription of the official XNA 4.0 documentation --- every
class, struct, enum, interface, and delegate across all 19 XNA namespaces, 544 types in
total, with their original property/method/constructor/field/event descriptions --- converted
from Microsoft's own \texttt{learn.microsoft.com} XNA~4.0 reference pages. Its own README
states its purpose directly: ``Auditing C\texttt{++} reimplementations of XNA~4.0 (comparing
what is implemented vs.\ what is missing).'' The coverage figures quoted in this book's
Chapter~\ref{ch:what-is-cna} and throughout Parts~I--IV (227 of 245 public FNA types present,
92.7\%) are exactly this kind of audit --- a diff against a ground truth, not an estimate.
Chapter~\ref{ch:xna4-spec-auditing} walks through using \texttt{xna4-spec} to run this kind of
audit yourself, worked through in full against a real class (\cnaclass{Ray}) rather than
described only in the abstract.

\section{cna-samples and cna-examples: two different kinds of proof}

Two sibling repositories exist purely to give CNA something real to run, and they answer two
different questions.

\texttt{cna-samples} is a C\texttt{++} port of the official Microsoft XNA Game Studio 4.0
sample collection (86 samples in the full upstream set), licensed Ms-PL to match its
Microsoft-derived origin. Its question is: \emph{does CNA run real, historical XNA sample
code?} --- and its current gap (23 of 86 samples blocked on the compiled \texttt{.fx}
bytecode gap from Chapter~\ref{ch:what-is-cna}) is one of the sharpest, most concrete
measurements of that gap that exists anywhere in the ecosystem.

\texttt{cna-examples} is original, MIT-licensed CNA-specific work --- explicitly \emph{not}
a port of Microsoft's sample collection, precisely so it can carry a permissive license
\texttt{cna-samples} cannot. Conceptually inspired by JavaFX's \texttt{Ensemble} sample
browser, it is a single cross-platform application (desktop, Web, mobile) that lets a reader
navigate Area~$\to$~Category~$\to$~Demo (Input, Audio, Devices, Net, Media, 2D/3D Graphics,
\ldots) and run a live demonstration of each area, all inside one binary. Its question is:
\emph{does CNA have a coherent, in-app way to show off everything it can do?} As of the
snapshot this book draws from it is an early skeleton --- the application boots and
navigates Area~$\to$~Category, but individual demo screens are not yet implemented, and only
the \texttt{EASYGL} backend is targeted so far.

\section{Repositories mentioned but out of this book's scope}

The \texttt{openeggbert} organization hosts a considerably larger set of repositories beyond
the six above --- additional CNA-adjacent projects (\texttt{cna-template}, \texttt{cna-extended},
\texttt{cna-craft}), other graphics libraries (\texttt{meta-gl}, \texttt{easy-3d}), actual
shipped games and tools built on this stack (\texttt{mesh-craft}, \texttt{galaxy-eggbert},
\texttt{free-eggbert}, \texttt{planetblupi}, \texttt{mobile-eggbert}), and a long tail of
per-project marketing/documentation sites (the \texttt{*.openeggbert.com} repositories). This
book's scope is deliberately the six repositories in the table above, because they are the
ones CNA's own build system, README, and \texttt{CLAUDE.md} porting rules directly name as
dependencies or verification partners. Appendix~\ref{ch:repo-map} gives a fuller repository
map for readers who want to go further.
